\section{Extremizers, two-point geometry and conservation}\label{j:core:section}
Throughout, $X,X_1,X_2,\ldots$ are iid, $\E X=0$, $\E X^2=1$,
$\rho=\E|X|^3<\infty$, and $S_n=\sum_{i=1}^nX_i$. Write
$\beta=\E[X|X|]$, $\gamma=\E X^3$, and $\phi=\Phi'$.
A threshold $H$ always refers to the raw sum; $z=H/\sqrt n$ is its
standardized value. Define the ratio and an auxiliary excess by
\[
 r(P,H)=\frac{\sqrt n}{\rho}\{\Pp(S_n\le H)-\Phi(z)\},\qquad
 a_c(P,H)=\sqrt n\{\Pp(S_n\le H)-\Phi(z)\}-c\rho.
\]
We maximize the ratio. The auxiliary quantity $a_c$ is used only to
state an identity for arbitrary coefficients and off-contact laws.

\subsection{Ratio attainment and contact, including atoms}
\begin{theorem}\label{j:core:attainment}
For every $n\ge2$, $C_n$ is attained as a positive closed-CDF ratio by
a compactly supported standardized law and a real raw threshold $H$.
At such a pair put $c=C_n$, $F=\Pp(S_n\le H)$,
$h(y)=\Pp(S_{n-1}\le H-y)$ and $\phi_z=\phi(z)$. Then
\begin{align}
 I_c(y)={}&n^{3/2}(h(y)-F)+ny\phi_z
 +\frac{H\phi_z+3c\rho}{2}(y^2-1)
 -c(|y|^3-\rho-3y\beta)\le0,\label{j:core:contact}\\
 Q(y)={}&F+\frac{H\phi_z+c\rho}{2n^{3/2}}
 -\frac{n\phi_z+3c\beta}{n^{3/2}}y
 -\frac{H\phi_z+3c\rho}{2n^{3/2}}y^2
 +\frac{c}{n^{3/2}}|y|^3.\label{j:core:Q}
\end{align}
Thus $\Delta=Q-h\ge0$ everywhere and $\Delta=0$ on the entire closed
support. No sign assumption on $H$ is needed.
\end{theorem}
\begin{proof}
Reflection takes a negative error to a positive closed error because
\[
 \Pp(-S_n\le-H)-\Phi(-z)
 =\Phi(z)-\Pp(S_n<H)\ge\Phi(z)-\Pp(S_n\le H).
\]
Hence $C_n=\sup r(P,H)$. Since $\rho\ge1$, $C_n\le\sqrt n$;
symmetric Bernoulli at $H=n$ shows $C_n>0$.
Choose $r_j\uparrow C_n$ with $r_j\ge C_n/2$. Then
$\rho_j\le2\sqrt n/C_n$, and the positive discrepancies are bounded
away from zero. Chebyshev on the negative tail and the normal upper
tail on the positive side bound $H_j$. Pass to $P_j\Rightarrow P$,
$H_j\to H$, $\rho_j\to b$. A common third-moment ceiling $M$ bounds
the first- and second-moment tails by $M/K^2$ and $M/K$, so the limit
retains mean zero and variance one. Lower semicontinuity gives
$\rho(P)\le b$. For the sum laws, the closed-set inequality, first at
$H+t$ and then as $t\downarrow0$, gives
\[
 \Pp(S_n\le H)-\Phi(H/\sqrt n)
 \ge\limsup_j\{\Pp(S_{n,j}\le H_j)-\Phi(H_j/\sqrt n)\}
 =C_n b/\sqrt n.
\]
Thus $r(P,H)\ge C_n b/\rho(P)\ge C_n$, proving attainment and
$\rho(P)=b$. This does not assume uniform integrability of third moments.

Contaminate the raw law by $(1-\epsilon)P+\epsilon\delta_y$, then
recenter and rescale. Its mean is $\epsilon y$, variance
$v_\epsilon=1+\epsilon(y^2-1)-\epsilon^2y^2$, and centered third
absolute moment has derivative $|y|^3-\rho-3y\beta$.
Dominated differentiation uses $3|y|(|x|+|y|)^2$.
Evaluate the standardized law at the image of the \emph{fixed raw event}
$S_n\le H$. Its probability is a polynomial in $\epsilon$, with
derivative $n(h(y)-F)$ even at an atom collision. The normal argument
$(H-n\epsilon y)/\sqrt{nv_\epsilon}$ has derivative
$-\sqrt n y-H(y^2-1)/(2\sqrt n)$. The ratio derivative is
$I_c(y)/\rho$, and maximality makes it nonpositive.
Furthermore $\int I_c\,dP=0$. Since the closed-event function $h$ is
upper semicontinuous, strict negativity at a support point would hold
on a neighborhood of positive mass. Thus $I_c=0$ on the closed support.
Finally $I_c(y)\le-c|y|^3+O(y^2+|y|+1)$, with constants depending on
this pair, excludes both sufficiently distant rays.
\end{proof}

At a ratio maximizer, put $D_0=F-\Phi(z)$ and $\lambda=D_0/\rho=C_n/\sqrt n$.
The majorant used in the finite calculations is simply
\begin{equation}\label{j:core:scaled-Q}
 nQ(y)=\lambda|y|^3-\frac{z\phi_z+3D_0}{2}y^2
 -(\sqrt n\phi_z+3\lambda\beta)y+nF+\frac{z\phi_z+D_0}{2}.
\end{equation}
Thus its equality on support is $n h=nQ$; it is not a majorant for a
different marginal. If $A\ge C_{n-1}$, this equality implies, at every
support point $t$,
\begin{equation}\label{j:core:predecessor}
 0\le\mathcal H_n(H,t)+C_n\{-|t|^3+\tfrac32\rho t^2+3\beta t+\rho e_n\}
       +\rho s_n n(A-C_n),
\end{equation}
where
\[
 s_n=\sqrt{n/(n-1)},\quad e_n=n(s_n-1)-\tfrac12,
\]
\[
 \mathcal H_n(H,t)=n^{3/2}\!\left\{\Phi\!\left(\frac{H-t}{\sqrt{n-1}}\right)
 -\Phi\!\left(\frac H{\sqrt n}\right)\right\}
 +nt\phi_z+\tfrac H2(t^2-1)\phi_z.
\]
Indeed substitute $h(t)\le\Phi((H-t)/\sqrt{n-1})+A\rho/\sqrt{n-1}$
in \eqref{j:core:contact} and use $\sqrt n(F-\Phi(z))=C_n\rho$.
The signed difference $A-C_n$ may be negative; replacing it by a
nonnegative upper bound gives a weaker permitted inequality.

For later induction, if $C_j\le c$ for $j<N$ and some $n\ge N$ has
$C_n\ge c$, the least such $n$ has an attained ratio pair with
$C_n\ge c$ and $C_j\le c$ for every $j<n$. This includes equality at the
selected index; it requires no supremum over varying sample sizes.

\subsection{One sign normalization and its moment constraints}
Put $p=\Pp(X>0)$, $q=1-p$, $v=pq$, and
$m=\E X_+=\E X_->0$. Assign a zero atom to the nonpositive class.
For $B=\ind_{\{X>0\}}$, let $V=\E[X\mid B]$, $W=X-V$, and
$\theta=\E W^2$. Then $V=-m/q$ or $m/p$, and
\begin{equation}\label{j:core:normalization}
 \ell=\frac m{pq}=\sqrt{\frac{1-\theta}{pq}},\quad
 U=W/\ell,\quad X=\ell(B+U-p)=\frac{B+U-p}{\sqrt{pq+s}},\quad
 s=\E U^2=\frac{pq\theta}{1-\theta}.
\end{equation}
Here $0<p,q<1$, $0\le\theta<1$ and $\E[U\mid B]=0$.
Write $U_i$ for the conditional laws, $w_i=\E U_i^2$,
$\mu_i=\E U_i^3$, and $\lambda=ns$. Thus $s=qw_0+pw_1$; if $s>0$,
$\alpha=pw_1/s\in[0,1]$ gives $w_0=s(1-\alpha)/q$, $w_1=s\alpha/p$.
If $s=0$, both conditional laws are zero and $X$ is standardized Bernoulli.
The exact common-marginal identities are
\begin{align}
 \rho/\ell^3&=v(1-2v)+3v(w_0+w_1)-q\mu_0+p\mu_1,\notag\\
 \gamma/\ell^3&=v(q-p)+3v(w_1-w_0)+q\mu_0+p\mu_1,\label{j:core:conditional-moments}\\
 \beta/\ell^2&=v(q-p)+pw_1-qw_0.\notag
\end{align}
They follow by expanding $\ell^3(p-U_0)^3$ and
$\ell^3(q+U_1)^3$ on the two sign classes and using conditional centering.
In particular the two cubic moments in the first two lines cannot be
selected independently. Higher moments require an established support
bound or another explicit integrability hypothesis.

\begin{lemma}\label{j:core:moments}
Every standardized law with finite third moment satisfies
\[
 \rho\ge\frac{\gamma^2+2}{\sqrt{\gamma^2+4}},\qquad
 |\beta|<1,\qquad \rho\ge\frac{1+\beta^2}{\sqrt{1-\beta^2}}.
\]
\end{lemma}
\begin{proof}
For the first inequality set $s_0=\sqrt{\gamma^2+4}$,
$a_0=(s_0-\gamma)/2$, $b_0=(s_0+\gamma)/2$ and
$c_0=\gamma(\gamma^2+6)/s_0^3$. The polynomial on each half-line
\[
 w(y)=|y|^3-c_0y^3-12y^2/s_0^3+6\gamma y/s_0^3+4/s_0^3
\]
has the factorizations, for $y,t\ge0$,
\[
 w(y)=\frac{2a_0^2}{s_0^3}(y-b_0)^2[2+(a_0+3b_0)y],\quad
 w(-t)=\frac{2b_0^2}{s_0^3}(t-a_0)^2[2+(3a_0+b_0)t].
\]
Taking expectations gives $0\le\E w=\rho-(\gamma^2+2)/s_0$.
For the second inequality let $a_\pm=\E X_\pm^2=(1\pm\beta)/2>0$.
Cauchy gives $m^2/a_++m^2/a_-\le1$, hence
$2m\le\sqrt{1-\beta^2}$, and
$m\rho\ge a_+^2+a_-^2=(1+\beta^2)/2$. These imply the result;
a zero atom only strengthens the first inequality.
\end{proof}

\subsection{A defect measuring departure from two points}
Let $K=\supp P$ and $\tau(t)=\E[X\ind_{\{X\ge t\}}]$.
Centering and signwise Tonelli give
\[
 \tau\ge0,\qquad \int\tau=1,\qquad 2\int|t|\tau(t)\,dt=\rho.
\]
On every bounded gap $(u,v)$ of $K$, $\tau$ has a constant interior value
$\tau_{u,v}$. Set
\begin{equation}\label{j:core:defect-definition}
 J=\sum_{(u,v)}\tau_{u,v}(v-u)^2,\qquad
 \sigma=\sqrt{10},\qquad D=(3+\sigma)\rho-\gamma-3J.
\end{equation}
The gaps are countable; the next bounds justify their sum even for
unbounded support.
\begin{theorem}\label{j:core:defect}
For every such law,
\[
 \gamma^2+4m^2\le\rho^2,\qquad J\le\rho+2m,\qquad
 D=(\sigma\rho-\gamma-6m)+3(\rho+2m-J)\ge0.
\]
Equality $D=0$ occurs exactly for
$X=(B-p_*)/\sqrt{p_*q_*}$, where
$p_*=(4-\sqrt{10})/2$ and $q_*=1-p_*$.
If $0\in K$, then $D\ge(\sigma-1)\rho$. Otherwise write the central
gap as $(-A_-,A_+)$ and define
\[
 T_\pm=\E X_\pm^3-mA_\pm^2,\quad
 R_\pm=\E X_\pm^2-mA_\pm,\quad J_o=J-m(A_-+A_+)^2.
\]
The quantities $T_\pm,R_\pm,J_o$ are nonnegative, and
\begin{align}
 J_o^\pm&\le T_\pm-2A_\pm R_\pm,\label{j:core:outer-gaps}\\
 D&=m\{\sqrt{\sigma+1}A_--\sqrt{\sigma-1}A_+\}^2
 +(\sigma-1)T_++(\sigma+1)T_-+3(T_++T_--J_o).\label{j:core:gap-defect}
\end{align}
All four terms in the last expression are nonnegative.
Moreover $1/(2\rho)\le m\le1/2$; if $D<2$, then
$T_++T_-\le D/2$ and $A_-+A_+\ge\sqrt{2-D}$.
\end{theorem}
\begin{proof}
Put $a_\pm=\E X_\pm^2$, $b_\pm=\E X_\pm^3$.
The signwise Cauchy inequalities give $a_\pm^2\le mb_\pm$ and
$m^2\le a_+a_-$. Thus $b_+b_-\ge m^2$, proving the circle inequality.
A gap lying on one side of zero satisfies
$(v-u)^2\le2\int_u^v|t|\,dt$. A central gap has weight $m$ and excess
$2mA_-A_+$ over this integral; since
$0\le\E[(X+A_-)(X-A_+)]=1-A_-A_+$, summation gives $J\le\rho+2m$.
Without a central gap it gives $J\le\rho$. Cauchy in $\R^2$ now gives
$\gamma+6m\le\sigma\rho$. Equality requires equality in both signwise
Cauchy bounds and in $m^2\le a_+a_-$: each sign is a single nonzero
atom and there is no zero atom. Writing its positive probability as $p$,
the remaining equality is $3(1-2p)=2p(1-p)$, whose unique solution
in $(0,1)$ is $p_*$. Conversely this two-point law gives equality.

For a positive outer gap,
$(v-u)^2\le2\int_u^v(t-A_+)\,dt$; summing yields
$J_o^+\le\E[X_+(X_+-A_+)^2]=T_+-2A_+R_+$.
The negative side is identical after reflection. Expansion using
$\sqrt{\sigma+1}\sqrt{\sigma-1}=3$ proves \eqref{j:core:gap-defect}.
Finally Cauchy gives $1\le2m\rho$ and $2m\le1$.
Since $\sigma-1>2$, $T_++T_-\le D/2$; hence
$m(A_-^2+A_+^2)=\rho-T_+-T_-\ge1-D/2$, which gives the last claim.
All gap summations use nonnegative Tonelli; endpoint atoms change no
Lebesgue integral.
\end{proof}
For the standardized Bernoulli law itself,
\[
 J=1/\sqrt{pq},\quad \rho_B=(1-2pq)/\sqrt{pq},\quad
 \gamma_B=(q-p)/\sqrt{pq},\quad
 D_B=(6+2\sigma)(p-p_*)^2/\sqrt{pq}.
\]
Thus the defect charges a quadratic displacement of $p$ as well as
conditional noise. The following bound charges them simultaneously.

\begin{lemma}[Joint defect budget]\label{j:core:joint-defect}
In a central gap put $\bar A_-=\ell p$, $\bar A_+=\ell q$,
$\Delta_\pm=\bar A_\pm-A_\pm\ge0$,
$a_0=\Delta_-/\ell$, $a_1=\Delta_+/\ell$, $d=a_0+a_1<1$, and
$w=1-\theta=pq\ell^2$. Then $U_0\le a_0$, $U_1\ge-a_1$, and
\begin{align}
 D\ge{}&\ell w(6+2\sigma)(p-p_*)^2
 +6\ell w\{d(1-d)+a_0a_1\}\notag\\
 &+\E[\{(\sigma-1)(X+2\bar A_+)+6A_+\}W^2;X>0]\notag\\
 &+\E[\{(\sigma+1)(-X+2\bar A_-)+6A_-\}W^2;X<0].\label{j:core:joint}
\end{align}
In particular, with $\theta_\pm=\E[W^2;\pm X>0]$,
\begin{align}
 D\ge{}&\frac{(6+2\sigma)(1-\theta)^{3/2}}{\sqrt{pq}}(p-p_*)^2
 +6m\{(A_-+A_+)(\Delta_-+\Delta_+)+\Delta_-\Delta_+\}\notag\\
 &+3(\sigma+1)A_+\theta_++3(\sigma+3)A_-\theta_-.\label{j:core:simultaneous}
\end{align}
\end{lemma}
\begin{proof}
Insert \eqref{j:core:outer-gaps} in the definition of $D$. The result is
\[
 D\ge(\sigma-1)M_{3,+}+(\sigma+1)M_{3,-}
 +6A_+M_{2,+}+6A_-M_{2,-}-6m(A_-^2+A_-A_++A_+^2),
\]
where $M_{j,\pm}=\E X_\pm^j$. Conditional centering gives
$M_{2,+}=m\bar A_++\E[W^2;X>0]$ and
$M_{3,+}=m\bar A_+^2+\E[(X+2\bar A_+)W^2;X>0]$,
and the reflected identities. Substitute $A_-=\ell(p-a_0)$,
$A_+=\ell(q-a_1)$; the deterministic part is precisely the first line
of \eqref{j:core:joint}. All kernels are nonnegative. Bounding them
below by $3(\sigma+1)A_+$ and $3(\sigma+3)A_-$ proves
\eqref{j:core:simultaneous}. Finite third moments suffice throughout.
\end{proof}

\subsection{Conservation: one identity, with and without contact}
For a compactly supported law retain $a=a_c$, $Q,h,\Delta$ from above,
and put $\nu=\mathcal L(H-S_{n-1})$. Over its bounded support gaps define
\[
 J_1=\sum\tau_{u,v}\!\int_u^v(t-u)t\,dt,\quad
 J_2=\sum\tau_{u,v}\!\int_u^v(t-u)t|t|\,dt,\quad
 L=\sum\tau_{u,v}\!\int_{[u,v)}(t-u)\,d\nu(t)\ge0.
\]
The half-open interval is essential: an atom at $v$ belongs to $h(v)$,
whereas an atom at $u$ receives zero weight. If $K\subset[-R,R]$, then
$J\le2R$, $|J_1|\le RJ/2$, $|J_2|\le R^2J/2$ and $L\le2R$.
Let $M_4=\E X^4$, $M_5=\E|X|^5$, and set
\[
 K_q=M_5/4-\rho M_4/2-3\beta\gamma/2+\rho/2,\quad
 K_1=\E[X|X|^3]-3\beta-3\rho\gamma/2,\quad
 \Omega=J_1+(M_4-3)/6.
\]
For $Z=S_n/\sqrt n$, a standard normal $Z_0$, and $q_z(x)=(x-z)|x-z|$, put
\[
 \varepsilon=\E q_z(Z)-\E q_z(Z_0)-2\gamma\phi_z/(3\sqrt n),\qquad
 E_n=-n^{3/2}\varepsilon/2,
\]
\begin{equation}\label{j:core:remainder}
 \mathcal R_c=c(3\beta J/2+3\rho J_1-3J_2-K_q)
 +H(cK_1+\phi_z\Omega)-H^2(c\rho+\gamma\phi_z/2)+E_n.
\end{equation}
Define the residuals
\[
 \xi=\E[X\Delta(X)],\qquad
 R_K=\int_K\tau(t)\Delta(t)\,dt
       +\sum_{(u,v)}\tau_{u,v}(v-u)\Delta(v).
\]
The integral on $K$ is a Lebesgue integral, not a $P$-integral. All terms
converge absolutely since $\Delta$ is bounded on the convex hull and
$\sum\tau_{u,v}(v-u)\le1$.

\begin{theorem}\label{j:core:conservation}
For every compactly supported standardized law and every $c>0$, $H\in\R$,
\begin{equation}\label{j:core:general-identity}
 (n+H^2)a+nc\rho-\frac{n\phi_z}{6}(\gamma+3J)+n^{3/2}L
 =\mathcal R_c+n^{3/2}(R_K-H\xi).
\end{equation}
At a ratio maximizer set $c=C_n$, so $a=0$. With
$e=C_n-\CE$ and $w=H^2/(2n)$, the residuals vanish and
\begin{equation}\label{j:core:positive-identity}
 \mathcal R_{C_n}=ne\rho+n\CE\rho(1-e^{-w})
       +n\phi(0)e^{-w}D/6+n^{3/2}L.
\end{equation}
All four terms are nonnegative when $C_n\ge\CE$, including equality.
\end{theorem}
\begin{proof}
Write $\mathcal A=\int\tau(Q-h)$. On a gap, $h(t)=h(v)+\nu([t,v))$.
Fubini and integration of $Q'$ therefore give
\begin{equation}\label{j:core:gap-integral}
 \mathcal A=R_K+\frac{\phi_zJ}{2\sqrt n}
 +\frac{3c\beta J/2+(H\phi_z+3c\rho)J_1-3cJ_2}{n^{3/2}}-L.
\end{equation}
For absolutely continuous $g$, signwise Fubini gives
$\int\tau g'=\E[Xg(X)]$. Thus
$\int t\tau=\gamma/2$, $\int t^2\tau=M_4/3$,
$\int|t|^3\tau=M_5/4$. Applying the same identity conditionally to
$g_s(x)=\min(s+x,H)$ and using exchangeability gives
\[
 \int\tau h=\frac1n\E[S_n\min(S_n,H)]
 =T,\qquad T=\E[(Z^2-zZ)\ind_{\{Z\le z\}}].
\]
The integrand at the threshold atom is zero. Consequently
\[
 \mathcal A=F-\frac{\gamma\phi_z}{2\sqrt n}
 +\frac{H\phi_z(3-M_4)/6+cK_q}{n^{3/2}}-T.
\]
If $M=\E[Z\ind_{\{Z\le z\}}]$, integrating $X(Q-h)$ against $P$ gives
\[
 M+\phi_z=-\gamma H\phi_z/(2n)+cK_1/n-\sqrt n\xi.
\]
Finally $q_z(x)=(x-z)^2-2(x-z)^2\ind_{\{x\le z\}}$, for both $Z$ and
$Z_0$, implies
\[
 T-\Phi(z)=-\varepsilon/2-\gamma\phi_z/(3\sqrt n)
               +z(M+\phi_z)-z^2(F-\Phi(z)).
\]
Substitute these two identities, use $\sqrt n(F-\Phi(z))=c\rho+a$,
and equate with \eqref{j:core:gap-integral}; collecting terms yields
\eqref{j:core:general-identity}. At contact $\Delta=0$ on the entire
support, so $R_K=\xi=0$. Substituting $a=0$, $c=C_n$ and
$\CE=(3+\sqrt{10})\phi(0)/6$ gives \eqref{j:core:positive-identity}.
\end{proof}
The off-contact identity permits checks on arbitrary bounded laws; it
is not a positivity statement there. In particular the support integral
in $R_K$ cannot be dropped for a continuous law. The proof uses contact
only when setting the two residuals to zero. Section~\ref{j:fourier:section}
proves uniform bounds for this exact $E_n$, including changing laws and
thresholds.

\begin{corollary}[The sign license]\label{j:core:license}
At a ratio maximizer with $C=C_n\ge\CE$, suppose a first pass proves
$\mathcal R_C\le T<nC\rho$. Then $\gamma+3J>0$ and
\begin{equation}\label{j:core:phase-price}
 n(C-\CE)\rho+n\phi(0)D/6+n^{3/2}L\le\mathcal R_C.
\end{equation}
\end{corollary}
\begin{proof}
The contact identity before splitting the defect gives
$\phi_z(\gamma+3J)/6=C\rho+\sqrt nL-\mathcal R_C/n>0$.
Thus $M=\CE\rho-\phi(0)D/6>0$. Regrouping
\eqref{j:core:positive-identity} writes its right side as the left side
of \eqref{j:core:phase-price} plus $n(1-e^{-w})M\ge0$.
\end{proof}
The initial bound is on the raw remainder. Only after it licenses the
positive phase term may that term improve a later resource bound.
No $D<2$ assumption or positive lower bound on $C_n-\CE$ was used.

For context, on a lattice with maximal span $h>0$, its gaps have lengths
$k_gh$, $k_g\ge1$. Since $\tau$ vanishes outside the convex hull and the
support has Lebesgue measure zero,
$\sum_g\tau_g k_gh=1$ and $J=\sum_g\tau_g k_g^2h^2\ge h$.
Reflection preserves $J,\rho,h$ and reverses $\gamma$, whence
$|\gamma|+3h\le(3+\sqrt{10})\rho$. Equality is precisely the zero-defect
Bernoulli law or its reflection. Thus the gap functional extends the
lattice span to arbitrary supports; no lattice hypothesis enters the proof.
